\subsection{A complete finite pentagon comparison in weight nine}

\begin{statement}{D.7}{Proved}{Compatible finite three- and four-strand shadows}
Write $P_n=PB_n$, $p=101$, and
\[
 K_n=\gamma_{11}P_n\,P_n^p\triangleleft B_n\qquad(n=2,3,4),
 \qquad K_F=\gamma_{11}F\,F^p,
\]
where a power subgroup is generated by all the indicated powers.
All these groups are discrete until their finite quotients are taken.
We use $F=\langle x=A_{12},y=A_{23}\rangle\le P_3$ and the same
braid generators as in D.4. Let $\mathcal G_9$ be the finite gentle
shadows at $K_3$ with $m=0\pmod p$ and
$f\in\gamma_9(F/K_F)$. Let $\mathcal P_9\subset\mathcal G_9$ be
those satisfying the pentagon in $P_4/K_4$.
The logarithm bijection makes both sets vector spaces over
$\mathbb F_{101}$ and gives the following dimensions:
\[
\begin{array}{c|rr}
 &\mathcal G_9&\mathcal P_9\\\hline
 \log f\text{ in degrees }9,10&37&2\\
 \text{degree-nine projection}&4&1\\
 \text{degree-nine metabelian projection}&2&1.
\end{array}
\tag{D.7.1}
\]
The last right-hand space is exactly $\mathbb F_{101}G_9[x,y]$.
The induced three-strand kernel associated to $K_4$ in the finite
parenthesized-braid definition is exactly $K_3$, so
$\mathcal P_9$ consists of actual finite pentagon shadows in that
definition, not merely solutions in an unrelated Lie target.
\end{statement}
\begin{proof}
\emph{The quotient maps.}
The subgroups $K_n$ are characteristic in $P_n$, hence normal in
$B_n$. The groups $P_n/K_n$ are finitely generated nilpotent groups
of bounded finite exponent, so are finite. Every strand insertion
or cabling homomorphism $P_3\to P_4$ sends $K_3$ into $K_4$,
since homomorphisms preserve the lower central series and power
subgroups. For the inclusion of the first three strands, the
strand-forgetting retraction $P_4\to P_3$ gives
$K_4\cap P_3=K_3$. Intersecting the five preimages used in
the definition of the induced three-strand kernel therefore gives
exactly $K_3$. The analogous two-strand assertion gives
$K_2=p\mathbb Z$ and $K_{\rm ord}=p$.
The direct product $P_3=F\times\langle C_3\rangle$ gives
$K_3\cap F=K_F$.
Thus reduction induces the specified maps of finite shadow sets,
and the continuous maps from the profinite braid groups to these
finite groups are their ordinary completion maps.

\emph{Faithful logarithmic coordinates.}
The class is ten, smaller than $p$. Lazard's Theorem II.4.6,
with its homomorphism assertion, identifies the finite groups
and their Lie rings by the mutually inverse truncated logarithm
and Campbell--Hausdorff operations \cite{DE_Lazard1954}.
For the free group factor this is the free nilpotent Lie algebra
of class ten over $\mathbb F_p$. This assertion can also be read
directly from Hall collection: the logarithm of a basic
commutator is its Hall Lie word plus terms of larger degree;
this triangular matrix has diagonal one, and all denominators
through degree ten are units modulo $p$.

For four strands we keep the actual pure Artin action in these
coordinates. Put $u_i=A_{i4}$, $a=A_{12}$, $b=A_{23}$.
Forgetting strand four splits $P_4$ as
$\langle u_1,u_2,u_3\rangle\rtimes P_3$.
The center is generated by
$C_4=C_3u_1u_2u_3$; eliminating $C_3$ in the quotient by
$C_4$ gives
\[
 P_4/\langle C_4\rangle
       =F(u_1,u_2,u_3)\rtimes F(a,b).
 \tag{D.7.2}
\]
These generators and their exact conjugation relations are
\cite[(2.1)--(2.3)]{DE_CohenFalkRandell2012}.
In particular the inverse actions of $a,b$ on the free factor are
\[
\begin{array}{c|ccc}
 &u_1&u_2&u_3\\\hline
 a^{-1}(-)a&u_1u_2u_1u_2^{-1}u_1^{-1}
                  &u_1u_2u_1^{-1}&u_3\\
 b^{-1}(-)b&u_1&u_2u_3u_2u_3^{-1}u_2^{-1}
                  &u_2u_3u_2^{-1}.
\end{array}
\tag{D.7.3}
\]
The actions are trivial on the free factor's abelianization.
Their Magnus expansions raise degree by at least one; a weight-$i$
commutator of them raises degree by at least $i$. Hence
$[\gamma_iF(a,b),\gamma_jF(u_1,u_2,u_3)]
\subset\gamma_{i+j}F(u_1,u_2,u_3)$.
It follows by both inclusions that the lower central series in
(D.7.2) is the product of the two free lower central series.
Hall collection and reduction modulo $p$ therefore give faithful
coordinates consisting of the two free Lie factors. The omitted
central factor has degree one. It intersects the derived subgroup
trivially, also after this reduction, and every pentagon factor
considered below lies in degree at least nine. Thus the computation
in (D.7.2) detects the complete pentagon in $P_4/K_4$.

\emph{The complete equations are linear in the retained stratum.}
Write $X=\log x$, $Y=\log y$ and
$\log f=h_9+h_{10}$, with $h_i$ homogeneous of degree $i$.
There are $56+99=155$ Lyndon coordinates.
The two simplified finite hexagons, with $z=y^{-1}x^{-1}$, are
\[
 f(x,y)f(y,x)=1,\qquad
 f(z,x)f(y,z)f(x,y)=1.
 \tag{D.7.4}
\]
Their equivalence to the original braid hexagons for $f\in F'$
is Dolgushev--Guynee, Proposition 3.4
\cite{DE_DolgushevGuynee2024}. Here $m=0$ and
\[
 \log z=-X-Y-\tfrac12[X,Y]\pmod{\text{degree }3}.
\]
Every product of two $h_i$ has degree at least eighteen.
Consequently (D.7.4), including all terms of degree ten, is
exactly a linear system. Substitution in the two-letter tensor
algebra produces a $3072$ by $155$ matrix of rank $118$.
Its kernel has dimension $37$; its degree-nine projection has
dimension four. The program also solves for, retains, and
substitutes a degree-ten correction for each of the four leading
basis vectors.

Here are the corresponding exact four-strand substitutions.
Put $U_i=\log u_i$, $A=\log a$, $B=\log b$, and
\[
 L=U_1+U_2+U_3+\tfrac12\sum_{i<j}[U_i,U_j]
             \pmod{\text{degree }3}.
\]
If the logarithm of an inverse Artin action in (D.7.3), applied
to $U_i$, is $U_i+v_{2i}+v_{3i}$ through degree three, then
its negative logarithm is the derivation
\[
 D(U_i)=-v_{2i}-v_{3i}
             +\tfrac12D_2(v_{2i}),\qquad D_2(U_i)=v_{2i}.
 \tag{D.7.5}
\]
This follows by expanding $-\log(1+T)=-T+T^2/2$ to the
required degree. The leading parts are
\[
\begin{aligned}
 D_A(U_1)&=[U_1,U_2],&D_A(U_2)&=[U_2,U_1],&D_A(U_3)&=0,\\
 D_B(U_1)&=0,&D_B(U_2)&=[U_2,U_3],&D_B(U_3)&=[U_3,U_2].
\end{aligned}
\]
Formula (D.7.5) supplies their degree-raising-two parts directly
from the literal group words. The five pairs of arguments of the
pentagon, in right-hand-side minus left-hand-side order, are
\[
\begin{array}{c|cc}
 &\text{first logarithm}&\text{second logarithm}\\\hline
 +&A&B+U_2+\tfrac12[U_2,U_3]\\
 +&-A-L&U_3\\
 -&B&U_3\\
 -&-B-L&U_2+U_3+\tfrac12[U_2,U_3]\\
 -&A&B.
\end{array}
\tag{D.7.6}
\]
For example $A_{13}A_{23}=a^{-1}C_3$ and
$A_{12}A_{13}=C_3b^{-1}$. The degree-two corrections involving
$A$ or $B$ and $L$ vanish because both leading derivations
annihilate $U_1+U_2+U_3$. Terms of degree three in the argument
logarithms, or derivation parts raising degree by three, cannot
affect a word of weight nine through degree ten. Therefore
(D.7.5)--(D.7.6) compute the complete group equation at this
quotient, without identifying the logarithms of the pure braid
generators with homogeneous infinitesimal generators.

The pentagon's free two-generator component cancels. Its free
three-generator component has $3^9$ and $3^{10}$ tensor
coefficients. The degree-nine equations reduce the $37$-dimensional
hexagon space to dimension $34$; all $59049$ degree-ten equations
then reduce it to dimension two. Projecting this kernel to degree
nine gives rank one, and applying the metabelian quotient gives
the asserted gamma line. All surviving basis vectors were
substituted into every hexagon and pentagon coefficient.

\emph{Generation.}
The map on $P_3/K_3$ induced by a retained pair sends
$x$ to $x$, $y$ to $f^{-1}yf$, and $C_3$ to $C_3$.
It is the identity on the Frattini quotient of this finite
$p$-group, and is therefore surjective. The two-strand map is
also surjective. Definition 2.9 and Proposition 2.10 of
Dolgushev--Le--Lorenz \cite{DE_DolgushevLeLorenz2024} now
give the generation condition on four strands for every retained
pentagon pair. The same three-strand argument gives the gentle
generation condition. This proves that the computed solutions
are the claimed finite shadows.

The executable sources are
\texttt{de\_gentle\_weight9.py} and
\texttt{de\_pentagon10\_finite.py}; their JSON outputs retain
the basis logarithms, and their compressed certificates retain
the full coefficient matrices. The latter run completed in
$501.307$ seconds. Its assertions verify every surviving
coefficient, the projection ranks, and equality with the gamma
line. No search over a subset of the $155$-coordinate stratum
is used.
\end{proof}

\begin{statement}{D.8}{Proved}{An integer pentagon certificate and the weight-nine comparison}
Let $k$ be a field of characteristic zero or characteristic at
least five. In the free Lie algebra $L_k(X,Y)$ put
\[
 V_9=\{h\in L_k(X,Y)_9:
 h(X,Y)+h(Y,X)=0,\quad
 [X,h(X,Y)]=[X+Y,h(-X-Y,Y)]\}.
\]
Then $\dim V_9=4$, its metabelian image has dimension two,
and the infinitesimal pentagon map $P:V_9\to(\mathfrak t_4)_9$
has rank three. Its one-dimensional kernel has metabelian image
$k\,G_9(X,Y)[X,Y]$. The complete finite projection in D.7 is this kernel
when $k=\mathbb F_{101}$.
For every $\ell\ge5$, applying the corresponding rational
leading-term map to the full profinite pentagon equation and
using D.20 gives precisely $\mathbb Q_\ell G_9(X,Y)[X,Y]$ as the
weight-nine metabelian image of $\widehat{\rm GT}$ and of
the Galois image.
\end{statement}
\begin{proof}
The following exact integer certificate also explains the
comparison with the finite calculation. Generate the $56$
Lyndon words of length nine lexicographically, and use standard
bracketing: split at the longest proper Lyndon suffix, then
bracket the two recursively bracketed factors. Expand in the
free associative algebra. The two equations defining $V_9$
give an integer $1536$ by $56$ matrix $C$. Its $52$ nonzero
Smith factors are all one. A four-column integer kernel basis
$H$ retained in the JSON certificate has Smith factors
$1,1,2,24$. Hence it is a basis of the kernel over every field
in the statement, not merely over $\mathbb Q$.

We evaluate the pentagon in the faithful Lie model
\[
 \mathfrak t_4/(\text{degree-one center})
       =L(U_1,U_2,U_3)\rtimes L(A,B),
\]
with the leading derivations displayed in D.7. It is the
associated graded of (D.7.2); the integral braid relations are
also recorded in \cite[(4.1)]{DE_CohenFalkRandell2012}.
For the five arguments take the degree-one parts of (D.7.6).
All eighteen infinitesimal braid relations are checked by the
program, as is the cancellation of the free $L(A,B)$ component.
The three tensor rows of $P(H)$ labelled
$111111223$, $111111232$, $111121123$ are
\[
 \begin{pmatrix}
 -6&-6&6&-5\\
 254&510&-446&-1291\\
 -150&-302&246&700
 \end{pmatrix}.
 \tag{D.8.1}
\]
The determinant of the first three columns is
$27648=2^{10}3^3$. The vector
\[
                 v=(-3299,1326,-1613,432)^{\mathsf T}
 \tag{D.8.2}
\]
annihilates all $19683$ tensor coefficients of $P(H)$, verified
with integer arithmetic. Its coordinates have greatest common
divisor one. Thus $P$ has rank exactly three and kernel $kv$
in each of the stated characteristics.

In the ordered monomial coordinates
$X^iY^{7-i}[X,Y]$ for $0\le i\le7$, the metabelian image of
$H$ is
\[
 \begin{pmatrix}
 0&0&0&-2\\
 0&0&0&-8\\
 2&2&-2&-17\\
 6&6&-6&-23\\
 6&6&-6&-23\\
 2&2&-2&-17\\
 0&0&0&-8\\
 0&0&0&-2
 \end{pmatrix}.
 \tag{D.8.3}
\]
Its nonzero Smith factors are $1,4$. Applied to (D.8.2), it
gives
\[
 -96(9,36,84,126,126,84,36,9)^{\mathsf T}=-96G_9.
\]
This proves all the dimension and coordinate assertions.
The programs \texttt{de\_pentagon9.py} and
\texttt{de\_pentagon9\_rational.py} construct the matrices
from the Lie relations; the second performs the integer and
rational calculations anew, including both Smith forms and
every zero coefficient.

For the final comparison let a full profinite element have
weight-nine leading action in the filtration of Part D.
Its $\ell$-adic multiplier is one. In the normalized gauge
the action fixes $x$ and its conjugator belongs to $F'$.
A lower-degree inner action fixing $x$ would have an inner
generator centralizing $X$. In degree at least two this is
zero in the free Lie algebra; the degree-one possibility is
excluded by the derived-subgroup normalization. The conjugator
therefore has no term below degree nine. Its leading word
belongs to $V_9$ by the peripheral and symmetry equations.

The insertion homomorphisms of pure braid groups preserve
their lower central series. Their associated graded maps
send each pure generator to the corresponding sum of
infinitesimal generators; this follows by abelianizing the
explicit product formulas in the five insertion maps.
In degree nine the product of pentagon factors becomes their
sum, since products of two degree-nine terms have degree
eighteen. The leading word consequently lies in the exact
kernel just computed over $\mathbb Q_\ell$. This is the
upper inclusion, with its source and target specified.
The actual Galois construction in D.20 supplies a nonzero
element in the same filtered full profinite source, giving
the reverse inclusion on rational spans. This derives R9 in
weight nine from the certified pentagon map.

Finally, other multipliers act nontrivially in degree one,
and lower nonzero leading words belong to lower filtration
pieces. They do not define a weight-nine leading class under
the map in Part D. This is why D.7 enumerates precisely the
stratum needed for this comparison. The quotient detects
an extra finite gentle direction, but the calculation does
not construct a compatible family in the full profinite
gentle group. The remaining existence question is D.2.
\end{proof}
